\scanpage{5}
Nous supposons maintenant $k$ de car. $0$, de sorte qu'on a la notion d'exponentielle d'un champ de vecteurs formel $\xi$, sur \deleted{\unsure{$X_n$}} un schéma. formel lisse $X$,
\[
(t,x)\longmapsto(\exp t\xi).x:\underset{\text{``axe temps''}}{\underline{X_1}}\times X\longrightarrow X
\]
\margin{Intégration d'équations différentielles}
g) Si $X$ est un $\mathcal A$-schéma formel (lisse) $\unsure{\xi|X}$ un\inserted{$\mathcal A_*$} champ de vecteurs, alors le résolvant $X_1\times X\unsure{-|\to}X$ est un $\mathcal A$-morphisme.

(Bien entendu, il suffirait de le \deleted{faire}\inserted{demander} pour $X$ de la forme $X_n\ldots$).

Ceci implique\inserted{et} (par l'\unsure{introd.} de la variable temps devenant variable spatiale), donc équivaut, à

g') Comme g), mais avec un ``champ de vecteurs variable avec le temps'' $\xi(t)$
(\diagram{\begin{tikzcd}
X_1\times X\arrow[r]\arrow[dr,dashed,"\operatorname{pr}_2"']&\underline{T(X)}\arrow[d,"p"']\\
&X
\end{tikzcd}}), on \unsure{réslout|résout} l'équation différentielle formelle
\margin{variété \deleted{\unsure{des tangentes}}\inserted{formelle} tangente de $X$}
\[
\frac{d}{dt}x(t)=\xi(t,x)
\]
La résolution $x=\varphi(t,x_0)$ correspond à un $\varphi:X_1\times X\unsure{-|\to}X$ qui est un $\mathcal A$-morphisme.
\mspage{4}
